\documentclass[11pt]{article}
\usepackage[margin=1in]{geometry}
\usepackage{amsmath,amssymb,amsthm,booktabs,array,longtable}
\usepackage[T1]{fontenc}
\usepackage{lmodern}
\usepackage[hidelinks,breaklinks]{hyperref}
\usepackage{xurl}
\newtheorem{theorem}{Theorem}
\newtheorem{proposition}[theorem]{Proposition}
\newtheorem{corollary}[theorem]{Corollary}
\DeclareMathOperator{\rad}{rad}
\title{Correction of OEIS A343881, comment 2, revision 16:\\
extensions of Kagey (2015) and Adenwalla (2025)}
\author{Haobo Ma\\\small ChronoAI Pte Ltd, Singapore
\and Wenlin Zhang\\\small National University of Singapore\\\small ChronoAI Pte Ltd, Singapore}
\date{November 9, 2026}
\begin{document}
\maketitle
\begin{abstract}
We correct the conjecture in comment 2 of Peter Kagey's OEIS A343881,
revision 16 (June 14, 2021), and extend results of Kagey (2015) and
Adenwalla (2025): the proposed value $72$ at $k=12$ is inadmissible
for every exponent $n\ge2$.
The exact prime-exponent minima are $27$, $18$, and $54$ at $p=2,3,5$,
respectively, and $144$ at every prime $p\ge7$.
More generally, writing $k=r^g$ with the prime-exponent vector of $r$
having greatest common divisor $1$, we prove that $T(p,k)=r^{g+1}$ for
all sufficiently large primes, with an explicit threshold.
This gives the precise range in which the proposed formula $k\rad(k)$
is correct. We also prove the entry's distinct-prime-row assertion and
first-column formula using the same valuation description. For the
cross-referenced sequence minimum in A343825, we prove eventual equality
with A071537 by an explicit determinant bound and reprove Adenwalla's
$k=12$ prime column: $20$ at $p=2$ and $18$ at every odd prime.
The valuation classes also answer the image questions in Kagey's
\emph{Open problem collection}, Problem 16: they characterize the cube-row
image and show that the two-value rows are injective exactly at prime exponents.
\end{abstract}

\section{The conjecture and its obstruction}
For integers $n\ge2$ and $k\ge1$, let $T(n,k)$ be the least integer $m>k$
such that
\begin{equation}\label{eq:works}
 k^x m^y=c^n,\qquad c\ge1,\quad 1\le x,y<n.
\end{equation}
This is the definition in the NAME of A343881~\cite{a343881}.
The minimum always exists: $m=k2^n$, $x=n-1$, $y=1$, and $c=2k$
satisfy~\eqref{eq:works}. For a prime $q$, write $v_q(a)$ for its
exponent in the positive integer $a$. Unique prime factorization gives
the equivalent condition
\begin{equation}\label{eq:valuation}
 n\mid x v_q(k)+y v_q(m)\quad\hbox{for every prime }q.
\end{equation}
Conversely, these divisibilities give the integer
$c=\prod_q q^{(xv_q(k)+yv_q(m))/n}$, so no direction of this equivalence
requires an additional condition.

The conjecture is in the second COMMENT of revision 16 of A343881:
\begin{quote}
``Conjecture: T(p,k) = A064549(k) for fixed k $>$ 1 and sufficiently large p.''
\end{quote}
It was added in revision 3 on May 2, 2021; revision 16 retains this wording.
Here $p$ is prime and A064549 is $k\rad(k)$, where
$\rad(k)=\prod_{q\mid k}q$~\cite{a064549}. The failure at $k=12$
has a two-coordinate proof.

\begin{proposition}\label{prop:obstruction}
For every $n\ge2$, no integers in~\eqref{eq:works} exist with
$k=12$ and $m=72$. In particular, $T(p,12)\ne72=12\rad(12)$ for
every prime $p$.
\end{proposition}
\begin{proof}
Since $12=2^2 3$ and $72=2^3 3^2$, condition~\eqref{eq:valuation}
would imply
\[
 n\mid 2x+3y,\qquad n\mid x+2y.
\]
Subtracting the first expression from twice the second gives $n\mid y$,
contrary to $1\le y<n$. The minimum exists by the construction above,
and a least admissible integer is itself admissible.
\end{proof}
There are primes above every bound, so Proposition~\ref{prop:obstruction}
negates the conjecture's eventual quantifier, not just a finite range.

\section{The corrected eventual formula}
Write $k=\prod_{q\mid k}q^{a_q}>1$ and put
\[
 g=\gcd_{q\mid k}a_q,\qquad r=\prod_{q\mid k}q^{a_q/g},\qquad
 M=r^{g+1},\qquad A=\max_{q\mid k}a_q,\qquad
 L=\lfloor\log_2 M\rfloor.
\]
Thus $k=r^g$ and the exponents of $r$ have greatest common divisor $1$.

\begin{theorem}\label{thm:correction}
For every prime $p>\max\{g+1,AL\}$,
\[
 T(p,k)=M=k r.
\]
Consequently, for fixed $k>1$, the conjectured eventual value $k\rad(k)$
is correct if and only if all the nonzero prime exponents of $k$ are equal.
\end{theorem}
\begin{proof}
The integer $M$ is admissible: take $x=p-(g+1)$, $y=g$, $c=k$.
The bound $p>g+1$ ensures $1\le x,y<p$, and
\[
 k^x M^y=r^{g(p-g-1)+(g+1)g}=r^{gp}=k^p.
\]
Suppose $k<m<M$ is admissible. Write $b_q=v_q(m)$, taking $a_q=0$
outside the primes of $k$. Since $m<M$, every $b_q\le L$.
For a prime $q\nmid k$, condition~\eqref{eq:valuation} gives $p\mid yb_q$.
As $1\le y<p$ and $p$ is prime, $p\mid b_q$; but $0\le b_q\le L<p$,
so $b_q=0$. Thus $m$ has no new prime factor.

For primes $q,s\mid k$, eliminate $x$ between their two congruences.
This gives
\[
 p\mid y(a_q b_s-a_s b_q),\qquad
 p\mid a_q b_s-a_s b_q.
\]
Both products in the difference lie in $[0,AL]$, so its absolute value
is at most $AL<p$. Therefore $a_q b_s=a_s b_q$ for every pair $q,s$.
The vectors $(b_q)$ and $(a_q/g)$ are rationally proportional.
Set $e_q=a_q/g$. B\'ezout's identity supplies integers $u_q$ with
$\sum_q u_qe_q=1$. If $b_q=t e_q$ with $t\in\mathbb Q$, then
$t=\sum_q u_qb_q\in\mathbb Z$. Since $m>1$, $t>0$.
Hence $m=r^t$, and $m>k=r^g$ forces $t\ge g+1$, a contradiction.

Finally, $kr=k\rad(k)$ holds exactly when $r=\rad(k)$.
By unique factorization this means $a_q/g=1$ for every $q\mid k$,
which is equivalent to equality of all the $a_q$.
\end{proof}
For $2\le k\le11$ all nonzero prime exponents are equal; $12$ is the
first exception. This explains why the entry's displayed columns $k\le10$
do not exhibit the obstruction. The general bound gives $p\ge17$
at $k=12$; the exact threshold is smaller.

\begin{corollary}\label{cor:column}
For primes $p$,
\[
 T(p,12)=
 \begin{cases}
 27,&p=2,\\
 18,&p=3,\\
 54,&p=5,\\
 144,&p\ge7.
 \end{cases}
\]
\end{corollary}
\begin{proof}
For $p=2$, necessarily $x=y=1$, and $12m$ is a square exactly when
$m=3s^2$ for a positive integer $s$. The least such $m>12$ is $27$.
For $p=3$, $12\cdot18=6^3$; none of $13,14,15,16,17$ works:
for $13,14,15,17$ a new prime occurs to exponent $1$, contradicting
$3\mid y$; for $16$ the $3$-valuation forces $3\mid x$.

For $p=5$, $12^2\cdot54=6^5$.
If $12<m<54$ worked, any new prime would occur to exponent at least $5$,
giving $m\ge5^5>54$. Both $2$ and $3$ must occur in $m$, since otherwise
one of $5\mid 2x$ or $5\mid x$ would hold. Thus $m=2^a3^b$, $a,b\ge1$.
The possibilities strictly between $12$ and $54$ are $18,24,36,48$;
the values of $2b-a$ are $3,-1,2,-2$, respectively.
Elimination in~\eqref{eq:valuation} requires $5\mid2b-a$, so none works.

For $p\ge7$, $m=144$ works with $x=p-2$, $y=1$, and $c=12$.
Suppose $12<m<144$ worked. A new prime factor would imply
$m\ge5^p>144$. Again both $2$ and $3$ must divide $m$, so
$m=2^a3^b$ with $a,b\ge1$. The bound on $m$ gives $a\le5$ and $b\le3$:
$2^6\cdot3>144$ and $2\cdot3^4>144$.
Thus $-3\le2b-a\le5$, whereas elimination requires $p\mid2b-a$.
It follows that $a=2b$, so $m=12^b$. No such integer is strictly
between $12$ and $144$.
\end{proof}

\section{Other statements in the posing entry}
The first COMMENT says, for prime $p$, that the
``row consists of distinct integers.''
Its $p=2$ case has Kagey's explicit 2015 proof in A072905~\cite{a072905};
A277781 records the $p=3$ assertion~\cite{a277781}.
The following argument treats every prime in the same way.

\begin{proposition}\label{prop:distinct}
For each prime $p$, the map $k\mapsto T(p,k)$ on positive integers
is injective.
\end{proposition}
\begin{proof}
Associate to $k$ its finitely supported valuation vector
$(v_q(k)\bmod p)_q$ over $\mathbb F_p$.
Two positive integers are equivalent when their vectors are both zero,
or are nonzero scalar multiples of each other.
This is an equivalence relation. By~\eqref{eq:valuation}, admissibility
for a prime exponent is exactly this relation: $x,y$ are nonzero modulo $p$,
and conversely a nonzero scalar ratio supplies such $x,y$.
The zero-vector case also works, with $x=y=1$.
Hence $T(p,k)$ is the next larger member of $k$'s equivalence class.
If $k_1<k_2$ had the same successor $m$, then $k_1,k_2,m$ would lie
in the same class, with $k_1<k_2<m$. This contradicts the minimality
of the successor of $k_1$.
\end{proof}

The same description gives an exact formula for any fixed prime.
For $1\le a<p$ put
\[
 C_a(k)=\prod_{q\mid k}q^{[a v_q(k)]_p},
 \qquad [b]_p\in\{0,\ldots,p-1\},\quad [b]_p\equiv b\pmod p.
\]
\begin{corollary}\label{cor:primeformula}
For every prime $p$ and positive integer $k$,
\[
 T(p,k)=\min_{1\le a<p} C_a(k)
 \left(\left\lfloor\sqrt[p]{k/C_a(k)}\right\rfloor+1\right)^p.
\]
\end{corollary}
\begin{proof}
An admissible valuation vector is $a(v_q(k))_q$ modulo $p$ for
some $a\ne0$. The positive integers with this residue vector are exactly
$C_a(k)t^p$, $t\ge1$: outside the support of $k$, all valuations are
multiples of $p$. The displayed expression is the least such integer
strictly exceeding $k$, minimized over the allowed residues.
\end{proof}
At $p=2$ this is Kagey's February 6, 2015 formula in A072905:
if $k=d t^2$ with $d$ squarefree, then $T(2,k)=d(t+1)^2$.
At $p=3$ it gives the minimum of the two cube-residue formulas in
A277781, FORMULA, namely A254767 and A277780. The latter's formula
$d(t+1)^3$ for $k=dt^3$, $d$ cubefree, is attributed there to
Hoft (November 2, 2016), with Corneth's proof dated November 3, 2016
\cite{a277780}.

\begin{corollary}\label{cor:squarerange}
The map $k\mapsto T(2,k)$ is a bijection from positive integers to
the integers that are not squarefree.
\end{corollary}
\begin{proof}
In each squarefree class, the formula sends $dt^2$ to $d(t+1)^2$.
Its range therefore consists exactly of $dt^2$ with $t\ge2$, each with
unique inverse $d(t-1)^2$. This reproves Kagey's June 22, 2015 result.
\end{proof}

The FORMULA of A343881 states \texttt{T(n,1) = 2\^{}A020639(n)};
A020639 supplies
the least-prime-factor notation~\cite{a020639}.
\begin{proposition}\label{prop:firstcolumn}
If $\ell$ is the least prime divisor of $n\ge2$, then $T(n,1)=2^\ell$.
\end{proposition}
\begin{proof}
If $m^y=c^n$ with $1\le y<n$, put $d=n/\gcd(n,y)>1$.
Then $d\mid v_q(m)$ for every prime $q$, so $m$ is a $d$th power
of an integer at least $2$. Since $d>1$ divides $n$, it satisfies
$d\ge\ell$, and $m\ge2^d\ge2^\ell$.
Equality is attained with $m=2^\ell$, $x=1$, $y=n/\ell$, and $c=2$.
\end{proof}

\begin{corollary}\label{cor:composite}
The row $k\mapsto T(n,k)$ is injective if and only if $n$ is prime.
For composite $n$, if $\ell$ is its least prime divisor, then
\[
 T(n,1)=T(n,2^{\ell-1})=2^\ell.
\]
\end{corollary}
\begin{proof}
The prime case is Proposition~\ref{prop:distinct}, and the first value
is Proposition~\ref{prop:firstcolumn}. Write $n=\ell h$, where $h\ge\ell$.
For $k=2^{\ell-1}$ and $m=2^\ell$, choose
$x=\ell$, $y=h-\ell+1$, and $c=2$. Both exponents lie in $[1,n-1]$,
and $(\ell-1)x+\ell y=n$.
If $k<m'<2^\ell$ were admissible, any prime $q\ne2$ dividing $m'$
would have $n/\gcd(n,y')>1$ dividing $v_q(m')$, where $y'$ is its
witness exponent. This divisor is at least $\ell$, forcing
$m'\ge3^\ell>2^\ell$. Thus $m'$ would be a power of $2$, but there is
no such power strictly between $2^{\ell-1}$ and $2^\ell$.
The two displayed starting arguments are distinct.
\end{proof}

\subsection*{The image questions in Problem 16}
Problem 16 of Kagey's \emph{Open problem collection} (August 4, 2025,
printed p.~16) asks how to ``characterize what integers are in A299117,''
the image of the cube row, and asks about squares, squarefree integers,
and general exponents~\cite{problems}. Its earlier January 16, 2020
version has the same question~\cite{problems2020}.

\begin{corollary}\label{cor:image}
For a prime $p$, let $J_p(m)=\min_{1\le a<p}C_a(m)$, with $C_a$
as in Corollary~\ref{cor:primeformula}. Then
\[
 m\in\operatorname{image}(T(p,\cdot))\quad\Longleftrightarrow\quad
 m>J_p(m).
\]
Every member of this image is nonsquarefree.
\end{corollary}
\begin{proof}
The positive members of the valuation class of $m$ are exactly the
integers $C_a(m)t^p$ with $1\le a<p$, $t\ge1$.
Their least member is $J_p(m)$. A member has a predecessor precisely
when it exceeds this least member; the greatest smaller class member
then maps to it under $T(p,\cdot)$.
If $m$ is squarefree, $C_a(m)=m^a$ for $1\le a<p$, so $J_p(m)=m$.
\end{proof}

\begin{corollary}\label{cor:cubeimage}
Write $m=A B^2t^3$, where $A,B$ are coprime squarefree positive integers
and $t\ge1$. Then
\[
 m\in\operatorname{A299117}
 \quad\Longleftrightarrow\quad t>1\ \hbox{or}\ A<B.
\]
In particular the cube-row image contains every square and every cube
greater than $1$, contains no squarefree integer, and excludes $1$.
For every prime $q$,
\[
 T(3,q)=\min\{q^2,8q\}.
\]
\end{corollary}
\begin{proof}
The two representatives are $C_1(m)=AB^2$ and $C_2(m)=A^2B$.
If $t>1$, then $C_1(m)<m$. If $t=1$, the image criterion reduces to
$A^2B<AB^2$, or $A<B$. A square with $t=1$ has $A=1$ and $B>1$
unless it is $1$; a cube greater than $1$ has $t>1$.
The squarefree exclusion follows from Corollary~\ref{cor:image}.
Finally the two classes for $k=q$ give the candidates $8q$ and $q^2$
in Corollary~\ref{cor:primeformula}.
\end{proof}
This answers Problem 16's image-characterization and squarefree questions.
The answer to its square question is yes for squares greater than $1$,
and no if $1$ is included. The same endpoint convention applies to its
cube note, whose formula $T(3,t^3)=(t+1)^3$ starts at $t\ge1$.
The 2020 version additionally states $a(p)=p^2$ for prime $p$ in its
note before Related question 1. The corrected value above gives
$a_3(11)=88<121$, while its conclusion that all prime squares are in
the image remains true. A299117's nonsquarefree assertion~\cite{a299117}
is proved by Corollary~\ref{cor:image}.
For Problem 16's generalization, the image criterion covers prime
exponents; the composite-exponent image characterization is not covered.

\section{The cross-referenced sequence minimum}
Write $S(n,k)$ for the minimum in A343825~\cite{a343825}: the least
endpoint of a nondecreasing sequence starting at $k$, with every multiplicity
less than $n$ and product an $n$th power. This is the function $g_n(k)$
in Adenwalla's Definition 2.2~\cite{adenwalla}. For $n\ge2$,
$S(n,0)=0$ and $S(n,1)=1$, using singleton sequences.
COMMENT 2 asserts \texttt{T(n,k) <= A343881(n,k)} in that entry's notation.

\begin{corollary}\label{cor:comparison}
For $n\ge2$ and $k\ge1$, $S(n,k)\le T(n,k)$.
For $n=2$ this inequality is strict for every $k\ge1$.
\end{corollary}
\begin{proof}
Use $x$ copies of $k$ and $y$ copies of an admissible $m>k$.
For strictness, if $k$ is square then $S(2,k)=k<T(2,k)$.
Otherwise write $k=dt^2$ with $d>1$ squarefree and $t\ge1$.
Put $s=dt+1$. The four distinct integers
\[
 dt^2<ts<dt(t+1)<(t+1)s
\]
have product $(dt^2(t+1)s)^2$, and
$(t+1)s<d(t+1)^2=T(2,k)$ since
$d(t+1)^2-(t+1)s=(d-1)(t+1)>0$.
This is the construction in Kagey and Rajesh's Theorem 32 (2024),
credited there to Rajesh's 2024 thesis~\cite{kageyrajesh}.
\end{proof}

Let $v(j)=(v_q(j))_q$. For a field $F$, define
\[
 E_F(k)=\min\{m\ge k: v(k)\in
 \operatorname{span}_F(v(k+1),\ldots,v(m))\}.
\]
Over $\mathbb F_p$, valuation vectors are reduced modulo $p$.
The empty span is $\{0\}$. For $k\ge1$,
\begin{equation}\label{eq:span}
 S(p,k)=E_{\mathbb F_p}(k),\qquad
 R(k):=\operatorname{A071537}(k)=E_{\mathbb Q}(k).
\end{equation}
Indeed, divide a nonzero coefficient of $v(k)$ out of a linear relation;
conversely, choose integer residues $0,\ldots,p-1$ of its coefficients
to obtain the required multiplicities. Over $\mathbb Q$, clear
denominators and use unique factorization to identify a valuation relation
with the signed-power product $1$ in A071537~\cite{a071537}.
These minima exist: for $k>1$, $v(k^2)=2v(k)$ proves existence over
$\mathbb Q$ and over every odd characteristic; characteristic $2$ follows
from $S(2,k)\le T(2,k)$. Also $E_F(1)=1$.

\begin{proposition}\label{prop:permutation}
For every prime $p$, $k\mapsto S(p,k)$ permutes the nonnegative
nonprime integers as $k$ ranges over the nonnegative integers.
Also $R$ is a bijection from the positive integers to the positive nonprimes.
\end{proposition}
\begin{proof}
The following elimination works over either field in~\eqref{eq:span}.
If $k<l$ had the same minimal endpoint $m$, both relations would have
nonzero coefficients at $m$ unless an endpoint equalled its start.
Eliminating $v(m)$ would express $v(k)$ using $v(k+1),\ldots,v(m-1)$,
a contradiction. If $l=m$ then $v(m)=0$, which can simply be omitted
from the relation for $k$, giving the same contradiction.
Thus the map is injective.

A prime endpoint $m$ is impossible: its own prime coordinate is zero
in every earlier vector, but is $1$ in $v(m)$.
For a composite $m$, $v(m)$ is in the span of the earlier prime vectors.
If $v(m)=0$, then $E_F(m)=m$. Otherwise choose the largest $k<m$ for
which $v(m)$ lies in the span of $v(k),\ldots,v(m-1)$.
The coefficient of $v(k)$ is nonzero, so $E_F(k)\le m$.
If $E_F(k)<m$, substituting its relation would remove $v(k)$, contrary
to maximality of $k$. Hence $E_F(k)=m$.
Finally $1$ maps to $1$, and the $S$ row also has $0$ mapping to $0$.
\end{proof}
For $S$, this reproves Adenwalla's Theorems 5.1 and 5.5 (2025).
The source assertion is A343825, COMMENT 1, ``a permutation of the
nonprime integers,'' for each prime row.
Its $p=2$ case is Reid's 1987 result, whose proof is reproduced by
Kagey and Rajesh (2024), Section 1~\cite{kageyrajesh}.
For $R$, it reproves Alekseyev's October 22, 2008 proof in A071537.

\begin{theorem}\label{thm:sequenceeventual}
Fix $k>1$, set $R=R(k)$, $d=\pi(R)$, $B=\lfloor\log_2 R\rfloor$,
and $H=d!B^d$. Then $S(p,k)=R$ for every prime $p>H$.
For $k=1$, equality holds for every prime.
\end{theorem}
\begin{proof}
For each $k\le m\le R$, form the integer valuation matrices with
columns $v(k),\ldots,v(m)$ and $v(k+1),\ldots,v(m)$, using the $d$
prime coordinates at most $R$. Each entry is at most $B$.
By the determinant expansion, every $s$-rowed minor has absolute value
at most $s!B^s\le H$. Consequently every nonzero minor remains nonzero
modulo $p>H$, and every zero minor remains zero. Both ranks are therefore
unchanged on reduction. Equality of these two ranks is equivalent to the
span condition in~\eqref{eq:span}. The first endpoint satisfying it is
thus the same over $\mathbb Q$ and $\mathbb F_p$, namely $R$.
\end{proof}
This proves COMMENT 3 of A343825, posed in revision 11 on May 2, 2021
and retained in revision 22 (June 14, 2021):
\begin{quote}
``Conjecture: T(p,k) = A071537(k) for fixed k and sufficiently large prime p.''
\end{quote}
Here $k\ge1$ is the domain of A071537. This is a different eventual
formula from COMMENT 2 of A343881: two allowed values and arbitrarily
many allowed values give different minima.

\begin{corollary}\label{cor:sequencecolumn}
For every prime $p$,
\[
 S(p,12)=\begin{cases}20,&p=2,\\18,&p\text{ odd}.\end{cases}
 \qquad R(12)=18.
\]
In particular $S(p,12)=T(p,12)$ exactly when $p=3$.
\end{corollary}
\begin{proof}
Before endpoint $18$, the coordinates at $13,7,5,17$ force coefficients
of $13,14,15,17$ to vanish. The only remaining later vector is
$v(16)$, whose $3$-coordinate is zero, so it cannot span $v(12)$.
The rational relation
\[
 8v(12)-3v(16)-4v(18)=0
\]
supplies endpoint $18$ over $\mathbb Q$ and over every odd characteristic,
since the coefficients $8$ and $4$ are then nonzero.
In characteristic $2$, the $3$-coordinate remains unavailable through
endpoint $19$: $v(16)=0$, $v(18)=v(2)$, and the prime $19$ cannot
participate. At endpoint $20$, $12\cdot15\cdot20=60^2$ supplies equality.
Compare with Corollary~\ref{cor:column} for the last assertion.
\end{proof}

This recovers Adenwalla's Appendix, item 8 (2025), for prime exponents.

\small
\begin{longtable}{@{}>{\raggedright\arraybackslash}p{0.43\linewidth}>{\raggedright\arraybackslash}p{0.49\linewidth}@{}}
\toprule
Source item & Coverage \\
\midrule
\endfirsthead
\toprule
Source item & Coverage (continued) \\
\midrule
\endhead
A343881, COMMENT 2 & Refuted and corrected: Proposition~\ref{prop:obstruction},
Theorem~\ref{thm:correction}. \\
A343881, COMMENT 1 & Proved: Proposition~\ref{prop:distinct}; includes
the A277781 distinctness assertion. \\
A343881, FORMULA & Proved: Proposition~\ref{prop:firstcolumn}. \\
A343881, example table & Finite check: rows $2\le n\le11$, $1\le k\le10$. \\
A072905, bijection COMMENT & Reproved; prior proof Kagey (2015):
Corollary~\ref{cor:squarerange}. \\
A072905, FORMULA; A277781, FORMULA & Reproved; prior formulas
Kagey (2015) and the cited entries: Corollary~\ref{cor:primeformula}. \\
A277780, formula COMMENT & Reproved; prior formula Hoft and proof
Corneth (2016): Corollary~\ref{cor:primeformula}, $p=3$ case. \\
A299117, nonsquarefree COMMENT & Proved:
Corollary~\ref{cor:image}. \\
Problem 16, main question & Answered: explicit cube-row image,
Corollary~\ref{cor:cubeimage}. \\
Problem 16, Related 1--2 & Answered: all squares $>1$, but not $1$;
no squarefree number. \\
Problem 16 (2020), note $a(p)=p^2$ & Refuted and corrected:
$a_3(q)=\min\{q^2,8q\}$, including $a_3(11)=88$. \\
Problem 16, general-exponent question & Prime images characterized;
composite images not covered. Injectivity iff prime is not posed;
proved in Corollary~\ref{cor:composite}. \\
A343825, COMMENT 1 & Reproved; prior proof Adenwalla (2025),
Theorems 5.1, 5.5: Proposition~\ref{prop:permutation}. \\
A343825, COMMENT 2 & Proved: Corollary~\ref{cor:comparison}.
Strict $n=2$ case reproves Rajesh and Kagey--Rajesh (2024). \\
A343825, COMMENT 3 & Proved with explicit bound:
Theorem~\ref{thm:sequenceeventual}. \\
Adenwalla, Appendix, item 8 & Reproved; prior proof Adenwalla (2025):
exact $k=12$ prime column in Corollary~\ref{cor:sequencecolumn}. \\
A071537, permutation COMMENT & Reproved; prior proof Alekseyev (2008):
Proposition~\ref{prop:permutation}. \\
A006255, other bounds and conjectures; A328045, composite row;
Adenwalla, Conjectures 5.6 and 6.10 & Not covered. \\
\bottomrule
\end{longtable}
\normalsize

\section{Verification coverage}
The declaration \texttt{result} in the pinned
\href{https://github.com/the-omega-institute/trureturing/blob/6974dba01daa1a25cc74b87a86c191e95fe2c681/D5/S3/Arith/Factorization/OeisA343881KageyPowerProducts.lean}{\texttt{OeisA343881KageyPowerProducts.lean}} source~\cite{lean}
proves the negation of A343881's eventual-equality conjecture.
Its definitions \texttt{Works}, \texttt{A064549}, and \texttt{claim}
match~\eqref{eq:works}, $k\rad(k)$, and the prime-threshold quantifiers,
respectively. The sibling declaration \texttt{result} in
\href{https://github.com/the-omega-institute/trureturing/blob/6974dba01daa1a25cc74b87a86c191e95fe2c681/D5/S3/Arith/Factorization/OeisA343825KageyPowerSequences.lean}{\texttt{OeisA343825KageyPowerSequences.lean}}~\cite{leansibling}
proves A343825's eventual equality with A071537 for every $k\ge1$.
Its \texttt{Valid}, \texttt{T}, \texttt{Relation}, and \texttt{A071537}
definitions give the sequence, endpoint minimum, signed multiplicative
relation, and relation minimum used above. Its conclusion is the eventual
equality, rather than the explicit determinant bound in
Theorem~\ref{thm:sequenceeventual}.
The corrected two-value formula, exact columns, row assertions, explicit
bounds, and first-column identity have the written proofs in this note.

The accompanying \texttt{verify.py} uses exact integer factorization and
divisibility. It checks the 100 displayed table entries, the $k=12$ column
at all primes below $48$, the obstruction for $2\le n\le60$, and the
corrected formula at $2\le k\le100$ for $p=31,37,41,43$.
An integer-root calculation independently checks admissibility at $k=12$,
$2\le n\le7$, $13\le m\le144$.
Further checks compare the fixed-prime formula with direct minimum searches,
the sequence minimum with exhaustive small multiplicity lists and the
source's displayed prime rows, all prime exponents below $100$ at $k=12$,
and the rational relation minimum with its $81$ published terms.
Rank computations also test the eventual sequence equality for
$2\le k\le30$ and $p=101,103,107,109$; these are finite instances,
not applications of the generally much larger bound $H$.
The image criterion is compared with all row predecessors through $m=500$
at $p=2,3,5,7$; its cube specialization is compared with the first $1000$
terms of A299117. Composite-row collisions are checked for every
composite $n\le50$.
Exact ranges and counts are in \texttt{scientific-checks.json}.
These finite checks are distinct from the all-input written proofs.

\begin{thebibliography}{99}\small\raggedright
\bibitem{a343881} P. Kagey, \emph{OEIS A343881}, created May 2, 2021,
revision 16, June 14, 2021; NAME, COMMENTS 1--2, FORMULA, and EXAMPLE.
\url{https://oeis.org/A343881}.
\bibitem{a064549} H. Bottomley, \emph{OEIS A064549}, created October 16, 2001,
revision 71, September 7, 2026; NAME and the example $a(12)=72$.
\url{https://oeis.org/A064549}.
\bibitem{a072905} B. Cloitre, \emph{OEIS A072905}, created August 10, 2002;
P. Kagey's COMMENT of June 22, 2015, proving bijectivity.
\url{https://oeis.org/A072905}.
\bibitem{a277781} P. Kagey, \emph{OEIS A277781}, created October 30, 2016;
NAME and the one-to-one COMMENT. \url{https://oeis.org/A277781}.
\bibitem{a020639} D. W. Wilson, \emph{OEIS A020639}, NAME (least prime divisor).
\url{https://oeis.org/A020639}.
\bibitem{a343825} P. Kagey, \emph{OEIS A343825}, created April 30, 2021,
revision 22, June 14, 2021; NAME and COMMENTS.
\url{https://oeis.org/A343825}.
\bibitem{a071537} A. Fink, \emph{OEIS A071537}, created May 30, 2002;
M. Alekseyev's proof and program, October 22, 2008; NAME and COMMENTS.
\url{https://oeis.org/A071537}.
\bibitem{a277780} P. Kagey, \emph{OEIS A277780}, created October 30, 2016;
H. F. W. Hoft's formula, November 2, 2016, and D. A. Corneth's proof,
November 3, 2016. \url{https://oeis.org/A277780}.
\bibitem{a299117} P. Kagey, \emph{OEIS A299117}, created February 2, 2018;
NAME and nonsquarefree COMMENT. \url{https://oeis.org/A299117}.
\bibitem{problems} P. Kagey, \emph{Open problem collection},
August 4, 2025, Problem 16, printed p.~16 (PDF p.~18).
\href{https://github.com/peterkagey/OpenishProblems/blob/4fcae7f356de0ee7f65db45ac07b97d4e1448a64/v2/mainfile.pdf}{Pinned collection PDF}.
\bibitem{problems2020} P. Kagey, \emph{Open problem collection},
January 16, 2020, Problem 16, printed p.~16 (PDF p.~17).
\url{https://www.gathering4gardner.org/g4g14gift/G4G14-PeterKagey-OpenishProblems.pdf}.
\bibitem{kageyrajesh} P. Kagey and K. Rajesh, \emph{On a Conjecture about
Ron Graham's Sequence}, arXiv:2410.04728v2 (October 11, 2024),
Sections 1--2 and Theorem 32.
\url{https://arxiv.org/abs/2410.04728v2}.
\bibitem{adenwalla} S. Adenwalla, \emph{On a Generalisation of a Function
of Ron Graham's}, arXiv:2504.19196v2 (May 8, 2025), Definition 2.2,
Theorems 5.1 and 5.5, and Appendix, item 8. First version April 27, 2025.
\url{https://arxiv.org/abs/2504.19196v2}.
\bibitem{lean}
\href{https://github.com/the-omega-institute/trureturing/blob/6974dba01daa1a25cc74b87a86c191e95fe2c681/D5/S3/Arith/Factorization/OeisA343881KageyPowerProducts.lean}{\emph{OeisA343881KageyPowerProducts.lean}}, declaration
\texttt{result}, pinned revision \texttt{6974dba}.
\bibitem{leansibling}
\href{https://github.com/the-omega-institute/trureturing/blob/6974dba01daa1a25cc74b87a86c191e95fe2c681/D5/S3/Arith/Factorization/OeisA343825KageyPowerSequences.lean}{\emph{OeisA343825KageyPowerSequences.lean}},
declaration \texttt{result}, pinned revision \texttt{6974dba}.
\end{thebibliography}
\noindent\small
License: \href{https://creativecommons.org/licenses/by/4.0/}{Creative Commons Attribution 4.0 International (CC BY 4.0)}.
\end{document}
